\begin{center}
\LARGE
Homotopy Methods for Solving Nonlinear Equations and
            their Applications in Circuit Simulation

\end{center}

\begin{center}
\Large
Ljiljana Trajkovic

\end{center}

\noindent
Date:  July 19th (Friday)\\
Time: 11:00-11:30\\

\begin{center}
\large
Abstract
\end{center}

Finding the dc operating points of transistor circuits is an important task in circuit simulation. The
problem is equivalent to solving sets of nonlinear algebraic equations describing transistor circuits.
Existing circuit simulators use Newton's method, or its variants, to achieve this task. Newton's
method is local and requires a good initial guess for convergence, while its variants globally
converge under restrictive conditions. Recent mathematical results guarantee the existence of
constructive, globally convergent homotopy methods for finding zeros of nonlinear maps with
probability one. We apply these results to the dc operating point problem by constructing various
homotopies to create a "simple" problem that is solved before proceeding with the continuation
process that will transform it into the initially stated "difficult" problem. These homotopies
proved successful in overcoming dc convergence problems often encountered in current circuit
simulators (ADVICE, TITAN). 

